% ======================================================================
%  Robótica Aérea — Sección: CONTROL
%  Control PI-D del cuadricóptero: deducción de las funciones de
%  transferencia en lazo cerrado y sintonización por asignación de polos.
% ======================================================================

\documentclass[11pt,a4paper]{article}

\usepackage[utf8]{inputenc}
\usepackage[T1]{fontenc}
\usepackage[spanish,es-noshorthands]{babel}
\usepackage[margin=2.5cm]{geometry}
\usepackage{amsmath,amssymb,amsfonts}
\usepackage{booktabs}
\usepackage{siunitx}
\usepackage{tikz}
\usetikzlibrary{arrows.meta,positioning,calc}
\usepackage[colorlinks=true,linkcolor=blue,citecolor=blue]{hyperref}

\sisetup{output-decimal-marker={.}, group-separator={\,}}

% ---- Estilos para el diagrama de bloques -----------------------------
\tikzset{
  block/.style = {draw, thick, rectangle, minimum height=2.6em, minimum width=2.6em},
  sum/.style   = {draw, thick, circle, inner sep=0pt, minimum size=6mm},
  input/.style = {coordinate},
  output/.style= {coordinate},
  punto/.style = {circle, fill, inner sep=1.2pt}
}

\newcommand{\Lap}{\mathcal{L}}

\title{\textbf{Control de actitud y posición de un cuadricóptero}\\[2mm]
\large Estructura PI--D, funciones de transferencia en lazo cerrado\\
y sintonización por asignación de polos}
\author{Sección de Control --- Robótica Aérea}
\date{\today}

\begin{document}
\maketitle

% ======================================================================
\section{Plantas en lazo abierto}
% ======================================================================

A partir del modelo linealizado alrededor del vuelo estacionario, los cuatro
canales quedan desacoplados y cada uno se reduce a un doble integrador:
\begin{equation}
\frac{\Phi(s)}{T_\phi(s)}=\frac{1/I_{xx}}{s^{2}},\qquad
\frac{\Theta(s)}{T_\theta(s)}=\frac{1/I_{yy}}{s^{2}},\qquad
\frac{\Psi(s)}{T_\psi(s)}=\frac{1/I_{zz}}{s^{2}},\qquad
\frac{P_N(s)}{F_t(s)}=\frac{1/m}{s^{2}} .
\label{eq:plantas}
\end{equation}

Las cuatro plantas comparten la misma estructura, por lo que se define la
\textbf{planta genérica}
\begin{equation}
G(s)=\frac{Y(s)}{U(s)}=\frac{1/J}{s^{2}},
\qquad
J \in \{\,I_{xx},\;I_{yy},\;I_{zz},\;m\,\},\quad
Y \in \{\,\Phi,\;\Theta,\;\Psi,\;P_N\,\},
\label{eq:generica}
\end{equation}
donde $J$ es el parámetro inercial del canal. Todo el desarrollo se realiza una
sola vez sobre \eqref{eq:generica} y luego se particulariza.

% ======================================================================
\section{Ley de control}
% ======================================================================

De acuerdo con el diagrama de bloques planteado, la ley de control es
\begin{equation}
u(t) = K_I \int_{0}^{t} e(\tau)\,d\tau \;-\; K_P\, y(t) \;-\; K_D\, \dot y(t),
\qquad e(t)=r(t)-y(t).
\label{eq:leycontrol}
\end{equation}

La acción integral actúa sobre el error $e=r-y$, mientras que las acciones
proporcional y derivativa actúan sobre la salida $y$. Esta configuración se
denomina \emph{PI--D}.

\begin{figure}[h!]
\centering
\begin{tikzpicture}[auto, node distance=1.5cm, >=Latex, thick]

  \node [input, name=input] {};
  \node [sum, right=1.4cm of input] (sum1) {};
  \node [block, right=1.2cm of sum1] (ki) {$K_I$};
  \node [block, right=1.2cm of ki] (int) {$\dfrac{1}{s}$};
  \node [sum, right=1.6cm of int] (sum2) {};
  \node [block, right=1.4cm of sum2, minimum width=3.2em] (plant) {$\dfrac{1/J}{s^{2}}$};
  \node [coordinate, right=1.0cm of plant] (fb) {};
  \node [output, right=1.0cm of fb] (output) {};

  \node [block, below=1.8cm of ki] (kp) {$K_P$};
  \node [block, below=1.8cm of kp] (kd) {$K_D$};
  \node [block, right=1.2cm of kd] (der) {$s$};

  % camino directo
  \draw [->] (input) -- node[pos=0.15,above] {$R(s)$} node[pos=0.85,above] {\small $+$} (sum1);
  \draw [->] (sum1) -- node[above] {\small $E(s)$} (ki);
  \draw [->] (ki) -- (int);
  \draw [->] (int) -- node[pos=0.85,above] {\small $+$} (sum2);
  \draw [->] (sum2) -- node[above] {$U(s)$} (plant);
  \draw [-]  (plant) -- (fb);
  \draw [->] (fb) -- node[above,pos=0.9] {$Y(s)$} (output);
  \node [punto] at (fb) {};

  % realimentación principal
  \draw [-]  (fb) -- ++(0,-5.9) coordinate (bottom);
  \draw [->] (bottom) -| node[pos=0.93,left] {\small $-$} (sum1);

  % ramas P y D
  \node [punto] at (sum1 |- kp) {};
  \node [punto] at (sum1 |- kd) {};
  \draw [->] (sum1 |- kp) -- (kp);
  \draw [->] (sum1 |- kd) -- (kd);
  \draw [->] (kd) -- (der);
  \draw [->] (kp.east) -| node[pos=0.93,left] {\small $-$} (sum2.south);
  \draw [->] (der.east) -| node[pos=0.93,right] {\small $-$} (sum2.south east);

\end{tikzpicture}
\caption{Estructura PI--D aplicada a la planta genérica $G(s)=\dfrac{1/J}{s^{2}}$.
Para el canal de alabeo, $J=I_{xx}$, $Y=\Phi$ y $R$ es la referencia de alabeo.}
\label{fig:bloques}
\end{figure}

% ======================================================================
\section{Función de transferencia en lazo cerrado}
\label{sec:lazocerrado}
% ======================================================================

\paragraph{Paso 1: transformada de Laplace de la ley de control.}
Aplicando $\Lap\{\cdot\}$ a \eqref{eq:leycontrol} con condiciones iniciales
nulas, y usando $\Lap\{\int_0^t e\} = E(s)/s$ y $\Lap\{\dot y\}=sY(s)$:
\begin{equation}
U(s)=\frac{K_I}{s}\,E(s)-K_P\,Y(s)-K_D\,s\,Y(s).
\end{equation}

\paragraph{Paso 2: sustitución del error.}
Con $E(s)=R(s)-Y(s)$,
\begin{equation}
U(s)=\frac{K_I}{s}\big[R(s)-Y(s)\big]-K_P\,Y(s)-K_D\,s\,Y(s).
\label{eq:paso2}
\end{equation}

\paragraph{Paso 3: ecuación de la planta.}
De \eqref{eq:generica},
\begin{equation}
Y(s)=\frac{1/J}{s^{2}}\,U(s)
\qquad\Longleftrightarrow\qquad
U(s)=J\,s^{2}\,Y(s).
\label{eq:paso3}
\end{equation}

\paragraph{Paso 4: igualación.}
Sustituyendo \eqref{eq:paso3} en \eqref{eq:paso2}:
\begin{equation}
J\,s^{2}\,Y(s)=\frac{K_I}{s}R(s)-\frac{K_I}{s}Y(s)-K_P\,Y(s)-K_D\,s\,Y(s).
\end{equation}

\paragraph{Paso 5: eliminación de denominadores.}
Multiplicando toda la ecuación por $s$:
\begin{equation}
J\,s^{3}\,Y(s)=K_I\,R(s)-K_I\,Y(s)-K_P\,s\,Y(s)-K_D\,s^{2}\,Y(s).
\end{equation}

\paragraph{Paso 6: agrupación de términos en $Y(s)$.}
\begin{equation}
Y(s)\Big[J\,s^{3}+K_D\,s^{2}+K_P\,s+K_I\Big]=K_I\,R(s).
\end{equation}

\paragraph{Paso 7: función de transferencia.}
\begin{equation}
\boxed{\;
\frac{Y(s)}{R(s)}=\frac{K_I}{J\,s^{3}+K_D\,s^{2}+K_P\,s+K_I}
=\frac{K_I/J}{s^{3}+\dfrac{K_D}{J}s^{2}+\dfrac{K_P}{J}s+\dfrac{K_I}{J}} \;}
\label{eq:lazocerrado}
\end{equation}

El lazo cerrado resulta de tercer orden: los dos polos de la planta más el polo
aportado por el integrador.

\subsection{Particularización a los cuatro canales}

Reemplazando $J$ y $Y$ en \eqref{eq:lazocerrado} se obtienen directamente las
cuatro funciones de transferencia.

\paragraph{Alabeo (\emph{roll}), $J=I_{xx}$:}
\begin{equation}
\boxed{\;\frac{\Phi(s)}{R_\phi(s)}
=\frac{K_{I\phi}}{I_{xx}s^{3}+K_{D\phi}s^{2}+K_{P\phi}s+K_{I\phi}}\;}
\end{equation}

\paragraph{Cabeceo (\emph{pitch}), $J=I_{yy}$:}
\begin{equation}
\boxed{\;\frac{\Theta(s)}{R_\theta(s)}
=\frac{K_{I\theta}}{I_{yy}s^{3}+K_{D\theta}s^{2}+K_{P\theta}s+K_{I\theta}}\;}
\end{equation}

\paragraph{Guiñada (\emph{yaw}), $J=I_{zz}$:}
\begin{equation}
\boxed{\;\frac{\Psi(s)}{R_\psi(s)}
=\frac{K_{I\psi}}{I_{zz}s^{3}+K_{D\psi}s^{2}+K_{P\psi}s+K_{I\psi}}\;}
\end{equation}

\paragraph{Posición norte, $J=m$:}
\begin{equation}
\boxed{\;\frac{P_N(s)}{R_{p}(s)}
=\frac{K_{Ip}}{m\,s^{3}+K_{Dp}s^{2}+K_{Pp}s+K_{Ip}}\;}
\end{equation}

Puesto que en un cuadricóptero simétrico $I_{xx}=I_{yy}$, los canales de alabeo
y cabeceo son idénticos y comparten la misma sintonización.

% ======================================================================
\section{Diseño por asignación de polos}
% ======================================================================

\subsection{Especificaciones de diseño}

Se adopta un amortiguamiento $\zeta=0.9$ para todos los canales, de modo que la
respuesta sea prácticamente sin oscilación. El tiempo de establecimiento se fija
según una jerarquía de lazos: el lazo interno de actitud debe ser el más rápido,
la guiñada intermedia y el lazo externo de posición el más lento.

\begin{center}
\begin{tabular}{lcc}
\toprule
Canal & $\zeta$ & $t_{ss}$ [s] \\
\midrule
Alabeo $\phi$ y cabeceo $\theta$ & 0.9 & 0.5 \\
Guiñada $\psi$                   & 0.9 & 1.0 \\
Posición $p_N$                   & 0.9 & 2.0 \\
\bottomrule
\end{tabular}
\end{center}

Con el criterio del $2\,\%$,
\begin{equation}
t_{ss}\approx\frac{4}{\zeta\omega_n}
\qquad\Longrightarrow\qquad
\omega_n=\frac{4}{\zeta\,t_{ss}} .
\label{eq:wn}
\end{equation}

El tercer polo se ubica en $p=5\zeta\omega_n$, suficientemente alejado del par
dominante para que su efecto transitorio sea despreciable.

\subsection{Polinomio característico deseado e igualación}

Los polos deseados corresponden a un par complejo conjugado dominante y un polo
real:
\begin{equation}
\Delta_d(s)=\big(s^{2}+2\zeta\omega_n s+\omega_n^{2}\big)\big(s+p\big).
\end{equation}

Desarrollando el producto:
\begin{align}
\Delta_d(s)&=s^{3}+2\zeta\omega_n s^{2}+\omega_n^{2}s
            +p\,s^{2}+2\zeta\omega_n p\,s+\omega_n^{2}p \nonumber\\[2pt]
&=s^{3}+\big(2\zeta\omega_n+p\big)s^{2}
       +\big(\omega_n^{2}+2\zeta\omega_n p\big)s
       +\omega_n^{2}p .
\label{eq:deltad}
\end{align}

El polinomio característico del lazo cerrado, tomado del denominador de
\eqref{eq:lazocerrado} en forma mónica, es
\begin{equation}
\Delta(s)=s^{3}+\frac{K_D}{J}s^{2}+\frac{K_P}{J}s+\frac{K_I}{J}.
\label{eq:delta}
\end{equation}

Igualando \eqref{eq:delta} con \eqref{eq:deltad} coeficiente a coeficiente:
\begin{equation}
\frac{K_D}{J}=2\zeta\omega_n+p,\qquad
\frac{K_P}{J}=\omega_n^{2}+2\zeta\omega_n p,\qquad
\frac{K_I}{J}=\omega_n^{2}p ,
\end{equation}
y despejando las ganancias:
\begin{equation}
\boxed{\;
K_D=J\big(2\zeta\omega_n+p\big),\qquad
K_P=J\big(\omega_n^{2}+2\zeta\omega_n p\big),\qquad
K_I=J\,\omega_n^{2}\,p \;}
\label{eq:ganancias}
\end{equation}

Las tres ganancias resultan proporcionales al parámetro inercial $J$ del canal.

\subsection{Aplicación a cada canal}

Evaluando \eqref{eq:wn} con $\zeta=0.9$ y $p=5\zeta\omega_n$:

\begin{center}
\begin{tabular}{lcccc}
\toprule
Canal & $t_{ss}$ [s] & $\omega_n$ [rad/s] & $\zeta\omega_n$ & $p$ [rad/s]\\
\midrule
$\phi,\ \theta$ & 0.5 & 8.889 & 8 & 40 \\
$\psi$          & 1.0 & 4.444 & 4 & 20 \\
$p_N$           & 2.0 & 2.222 & 2 & 10 \\
\bottomrule
\end{tabular}
\end{center}

Sustituyendo estos valores en \eqref{eq:ganancias} se obtienen las ganancias en
forma simbólica:

\begin{center}
\begin{tabular}{lccc}
\toprule
Canal & $K_P$ & $K_I$ & $K_D$ \\
\midrule
$\phi$   & $719.01\,I_{xx}$ & $3160.49\,I_{xx}$ & $56\,I_{xx}$ \\
$\theta$ & $719.01\,I_{yy}$ & $3160.49\,I_{yy}$ & $56\,I_{yy}$ \\
$\psi$   & $179.75\,I_{zz}$ & $\phantom{0}395.06\,I_{zz}$ & $28\,I_{zz}$ \\
$p_N$    & $\phantom{0}44.94\,m$ & $\phantom{00}49.38\,m$ & $14\,m$ \\
\bottomrule
\end{tabular}
\end{center}

\subsection{Valores numéricos}

Empleando los momentos de inercia deducidos en la sección de modelado,
\begin{equation}
I_{xx}=I_{yy}=\frac{2MR^{2}}{5}+2\ell^{2}m_r,
\qquad
I_{zz}=\frac{2MR^{2}}{5}+4\ell^{2}m_r,
\qquad
m=M+4m_r ,
\end{equation}
con $M=\SI{1.0}{\kilogram}$, $R=\SI{0.05}{\meter}$, $\ell=\SI{0.225}{\meter}$ y
$m_r=\SI{0.06}{\kilogram}$:
\begin{equation}
I_{xx}=I_{yy}=\SI{7.075e-3}{\kilogram\meter\squared},\qquad
I_{zz}=\SI{1.315e-2}{\kilogram\meter\squared},\qquad
m=\SI{1.24}{\kilogram}.
\end{equation}

\begin{center}
\begin{tabular}{lcccc}
\toprule
Canal & $J$ & $K_P$ & $K_I$ & $K_D$ \\
\midrule
$\phi$   & \num{7.075e-3} & \num{5.0870}  & \num{22.3605} & \num{0.39620} \\
$\theta$ & \num{7.075e-3} & \num{5.0870}  & \num{22.3605} & \num{0.39620} \\
$\psi$   & \num{1.315e-2} & \num{2.3638}  & \num{5.1951}  & \num{0.36820} \\
$p_N$    & \num{1.24}     & \num{55.7235} & \num{61.2346} & \num{17.3600} \\
\bottomrule
\end{tabular}
\end{center}

Las funciones de transferencia en lazo cerrado quedan entonces
\begin{align}
\frac{\Phi(s)}{R_\phi(s)}=\frac{\Theta(s)}{R_\theta(s)}
&=\frac{3160.5}{s^{3}+56\,s^{2}+719.0\,s+3160.5},\\[6pt]
\frac{\Psi(s)}{R_\psi(s)}
&=\frac{395.06}{s^{3}+28\,s^{2}+179.75\,s+395.06},\\[6pt]
\frac{P_N(s)}{R_p(s)}
&=\frac{49.383}{s^{3}+14\,s^{2}+44.938\,s+49.383},
\end{align}
cuyos polos son, por construcción, los asignados:

\begin{center}
\begin{tabular}{lcc}
\toprule
Canal & Par dominante & Polo real \\
\midrule
$\phi,\ \theta$ & $-8 \pm j3.874$ & $-40$ \\
$\psi$          & $-4 \pm j1.937$ & $-20$ \\
$p_N$           & $-2 \pm j0.968$ & $-10$ \\
\bottomrule
\end{tabular}
\end{center}

\noindent
donde la parte imaginaria corresponde a
$\omega_d=\omega_n\sqrt{1-\zeta^{2}}$.

\end{document}
